\section{Discussion and limitations}

For the binary observational class in \cref{obj:def:complete-model}, \cref{obj:thm:matched-risk-frontier,obj:thm:sharp-interval-frontier} give the common order \(r(n,\epsilon)\) for private point estimation and honest expected interval length. The lower bounds range over every procedure admitted by the respective decision criteria, while the pair-ratio release supplies attainment.

The shared occupancy weighting is the main reusable feature of the upper-bound construction. It gives the population numerator and denominator common density-dependent occupancy weights, while the occupancy scale \(\mathcal A=nh\min\{1,n\delta\}\) measures the availability of usable pairs in localized cells. Nuisance roughness determines the required cell resolution, and privacy noise is then amplified by the same occupancy shortage. This construction shows how density robustness, rough nuisances, and privacy interact through the occupancy scale.

The two private scales have different origins. The term \((n^2\epsilon)^{-1/7}\) reflects privacy in the sparse-pair regime, while \((n\epsilon)^{-1/2}\) reflects privacy once cells are densely occupied. Their transition to saturation yields the consistency condition \(n\epsilon_n\to+\infty\), and the regime comparison characterizes optimal dependence on sample size and privacy budget within the specified model.

Honest uncertainty quantification shares the estimation benchmark through complementary upper and lower arguments. The upper bound converts the uniform squared-error envelope into an interval radius. For the converse, transferred coverage forces a connected interval to contain two separated targets with positive probability under a common release distribution, and connectedness converts that separation into expected length. This simultaneous-containment argument establishes the interval frontier directly. Privacy is uniform over adjacent datasets, while error and coverage are uniform over laws in the causal class and include release randomization.

\subsection{Limitations and future work}

The certified interval constants are extremely conservative. As shown after \cref{obj:def:public-tuning}, the local-branch radius already exceeds two for every \(2\le n\le2^{40}\), so the reported interval is deterministically \([-1,1]\) throughout that range; the fallback branch is also full-range. The interval frontier should therefore be read as an asymptotic order result and a finite-sample coverage certificate, rather than as a claim of useful finite-sample width at these constants.

The regime identities characterize rates and do not by themselves prescribe finite-sample sample-size or privacy-budget choices.

The release is specified using real-valued arithmetic and ideal Laplace random variables. A represented-real implementation requires explicit treatment of numerical rounding, noise generation, denominator stabilization, and interval endpoints. Software privacy certification is a separate task: it must establish the privacy guarantee for the implemented output distribution and account for the numerical operations used by the program.

The minimax comparisons range over pure-private procedures with budgets in \((0,1]\). Establishing the unrestricted nonprivate minimax risk is a separate comparison. In particular, the benchmark at \(\epsilon=1\) describes the stated private decision problem; identifying its relation to an unrestricted nonprivate optimum requires an appropriate converse and matching construction for that larger procedure class.

The model uses binary treatment and outcomes, a scalar covariate, independent records, and fixed density, overlap, and smoothness restrictions. Broader outcome distributions, higher-dimensional covariates, and dependent or alternative sampling designs call for corresponding identification conditions, sensitivity bounds, and lower-bound experiments. Allowing overlap to vary with sample size or changing the nuisance and contrast smoothness classes would likewise require a new analysis of the bias, occupancy, and privacy balance. Adaptation to unknown smoothness introduces the additional question of selecting resolution while preserving privacy and honest coverage.

Finally, the results evaluate worst-case decision criteria and supply a concrete procedure attaining their optimal orders. Evaluation of particular causal-effect estimators remains a useful direction, including comparisons of finite-sample error, coverage, interval length, and computational cost under specified data-generating laws. Such evaluations would clarify how the explicit construction and its conservative error envelope perform at sample sizes and privacy budgets relevant to empirical applications.
